%% ============================================================================
%% CAPÍTULO 1 - EL PROBLEMA
%% ============================================================================

\chapter{EL PROBLEMA}

\section{Planteamiento del problema}

Los robots colaborativos (cobots) son una tecnología emergente que está revolucionando la industria 4.0, gracias a su capacidad de operar de manera segura y eficiente junto a los seres humanos. Esta característica clave ha permitido no solo avanzar hacia la automatización de procesos, sino también consolidar el concepto de la industria 5.0, donde la colaboración humano-máquina es más estrecha y fluida.

En este contexto, el diseño y desarrollo de un cobot dentro de la Universidad Tecnológica del Uruguay (UTEC) no solo impulsará la adopción de esta tecnología, sino que también generará beneficios significativos. Entre estos se incluyen la creación de nuevos proyectos de investigación en robótica colaborativa y el fortalecimiento del aprendizaje práctico de los estudiantes de ingeniería en mecatrónica.

Los cobots se destacan por su flexibilidad, tamaño reducido y fácil integración en diversos entornos laborales. Sin embargo, diseñar y construir un cobot presenta varios desafíos técnicos. Estos incluyen la selección adecuada de componentes mecánicos, eléctricos y electrónicos, así como la integración efectiva de sistemas de control. En este proyecto se plantea el uso del motor DRS 0602, conocido por su precisión y rápida capacidad de respuesta, como solución para la propulsión del cobot. Además, se utilizará el gripper piCOBOT, que está diseñado específicamente para aplicaciones colaborativas y ofrece una solución eficaz para la manipulación de objetos. La integración de estos componentes requiere un análisis en profundidad de todos los aspectos mecánicos y electrónicos, y un diseño cuidadoso de la estructura y el sistema de control del cobot.

Además de ser una justificación técnica y académica, este proyecto representa una valiosa oportunidad para nosotros como estudiantes de ingeniería en mecatrónica. Nos permite aplicar y desarrollar los conocimientos adquiridos a lo largo de nuestra formación académica en un entorno práctico. Al enfrentar estos retos, podremos adquirir competencias clave en robótica colaborativa y tecnologías emergentes, contribuyendo así al progreso del sector y a la transición de sistemas de producción actuales hacia la industria 4.0.

\section{Objetivos}

\subsection{Objetivo General}

Diseñar y construir un cobot (robot colaborativo) utilizando el motor DRS 0602 y el gripper piCOBOT.

\subsection{Objetivos Específicos}

\begin{enumerate}
  \item Analizar los aspectos mecánicos, eléctricos y electrónicos relevantes para el diseño y funcionamiento del robot.
  \item Diseñar la estructura del robot y un sistema básico de control integrando un sistema embebido, considerando aspectos de robótica colaborativa.
  \item Construir la estructura y la electrónica, integrando todos los componentes necesarios para el funcionamiento del robot colaborativo.
  \item Verificar el funcionamiento de la estructura y el sistema de control para el robot colaborativo con los motores DRS 0602 y el gripper piCOBOT.
\end{enumerate}

\section{Alcance, limitaciones y beneficiarios}

\subsection{Alcance}

% Completar según el documento original.
Describa aquí el alcance del proyecto.

\subsection{Limitaciones}

% Completar según el documento original.
Describa aquí las limitaciones encontradas.
